\section{Gaussian directions: the probabilistic proof}
\label{app:directions}

\begin{proof}[Proof of \Cref{lem:directions}]
Choose all the $g_{s,i}$ independently with the standard Gaussian distribution
on $\R^m$.  We show that the required events hold simultaneously with
probability tending to one as $m\to\infty$.

We first record the elementary estimates used below.  For a standard Gaussian
$G$ in $\R^m$,
\begin{equation}\label{dir:gaussian-identities}
 \E\cos\langle G,w\rangle=e^{-|w|^2/2},
 \qquad
 \E\bigl[G\sin\langle G,w\rangle\bigr]=w e^{-|w|^2/2}.
\end{equation}
Indeed, rotation reduces the first identity to the one-dimensional cosine
transform.  Gaussian integration by parts gives the differential equation
$\phi'(t)=-t\phi(t)$ for that transform, with $\phi(0)=1$.  Differentiating
the first identity with respect to $w$ gives the second; the finite first
Gaussian moment justifies this differentiation.  Also, exponential Markov
inequalities applied to
$\E e^{t|G|^2}=(1-2t)^{-m/2}$, for $t<1/2$, give
for an absolute $c>0$
\begin{equation}\label{dir:gaussian-norm-tail}
 \Pr\left\{|G|\notin[\sqrt m/2,2\sqrt m]\right\}
 \le 2e^{-c m}.
\end{equation}

We will use the following bounded-variable estimate, a form of
Hoeffding's inequality~\cite{Hoe63}. If $X_1,\ldots,X_n$
are independent real random variables with $|X_j|\le B$, then, for
$0<u\le B$,
\begin{equation}\label{dir:bounded-concentration}
 \Pr\left\{\left|\frac1n\sum_{j=1}^n(X_j-\E X_j)\right|>u\right\}
 \le 2\exp\left(-\frac{c n u^2}{B^2}\right).
\end{equation}
For completeness, Taylor expansion, the vanishing first centered moment,
and $|X_j-\E X_j|\le2B$ give
$\E e^{t(X_j-\E X_j)}\le e^{C t^2B^2}$ when $|t|\le1/(2B)$.
Independence and exponential Markov, with $t$ a sufficiently small constant
multiple of $u/B^2$, give both tails in
\eqref{dir:bounded-concentration}.

To obtain \eqref{dir:gaussian-approximation}, truncate each Gaussian summand
on the event $|G|>2\sqrt m$.  By Cauchy--Schwarz and
\eqref{dir:gaussian-norm-tail}, uniformly in $w$,
\begin{equation}\label{dir:truncation-bias}
 \left|\E\bigl[G\sin\langle G,w\rangle
                   \mathbf 1_{\{|G|\le2\sqrt m\}}\bigr]
               -w e^{-|w|^2/2}\right|
 \le \E\bigl[|G|\mathbf 1_{\{|G|>2\sqrt m\}}\bigr]
 \le e^{-c' m},
\end{equation}
after decreasing the positive constant $c'$ and taking $m$ large.
For a fixed unit vector $v$ and fixed $w$, the scalar truncated summand
\[
 \langle v,G\rangle\sin\langle G,w\rangle
       \mathbf 1_{\{|G|\le2\sqrt m\}}
\]
has absolute value at most $2\sqrt m$.  Thus
\eqref{dir:bounded-concentration}, with $n=N$ and
$u=1/(2\sqrt m)$, bounds its empirical-mean deviation probability by
$2e^{-cN/m^2}$.

Choose a $1/8$-net $\mathcal V$ of the unit sphere and an $m^{-3}$-net
$\mathcal W$ of the ball of radius $m^{42}$.  The usual disjoint-ball packing
argument gives
\[
 |\mathcal V|\le17^m,
 \qquad
 |\mathcal W|\le(1+2m^{45})^m
                 =\exp(O(m\log m)).
\]
Union bound the preceding scalar deviations over $s$, $v\in\mathcal V$,
and $w\in\mathcal W$.  The failure probability tends to zero, since
$N/m^2=m^4$.  On the resulting event, for every $s$ and $w\in\mathcal W$,
the vector deviation from its truncated expectation is at most
$4/(7\sqrt m)$: bounds on inner products with a $1/8$-net control the
Euclidean norm with factor $8/7$.

Each truncated empirical vector function is $4m$-Lipschitz in $w$, because
the operator norm of its derivative is at most the average of the
truncated $|g_{s,i}|^2$.  The function $w\mapsto we^{-|w|^2/2}$ is
$2$-Lipschitz.  Consequently the bias bound
\eqref{dir:truncation-bias} and interpolation from $\mathcal W$ give,
uniformly on the whole ball, an error at most
\[
 \frac{4}{7\sqrt m}+e^{-c'm}+(4m+2)m^{-3}
 \le \frac2{\sqrt m}
\]
for all sufficiently large $m$.  This proves the desired approximation
for the truncated samples.

The covariance estimate follows from the same argument with quadratic
forms.  For a fixed unit vector $v$, the variable
$\langle v,G\rangle^2\mathbf 1_{\{|G|\le2\sqrt m\}}$ is bounded by $4m$,
and its expectation differs from $1$ by at most
\[
 \bigl(\E\langle v,G\rangle^4\bigr)^{1/2}
 \Pr\{|G|>2\sqrt m\}^{1/2}
 \le e^{-c'm},
\]
where constants can again be adjusted, since
$\E\langle v,G\rangle^4=3$.  Apply
\eqref{dir:bounded-concentration} with $u=1/8$, and union bound over $s$
and $v\in\mathcal V$.  With probability tending to one, every truncated
empirical covariance matrix $M_s$ satisfies
\[
 \|M_s-I_m\|_{\op}
 \le\frac{1}{1-2/8}\left(\frac18+e^{-c'm}\right)
 <\frac12.
\]
Here we used the elementary symmetric-matrix net estimate
$\|A\|_{\op}\le(1-2\rho)^{-1}\max_{v\in\mathcal V}|\langle Av,v\rangle|$
for a sphere net of radius $\rho<1/2$.

By \eqref{dir:gaussian-norm-tail}, all $SN=m^9$ original Gaussian samples
satisfy \eqref{dir:norms} with probability tending to one.  On this event
the truncations change no sample, so the preceding conclusions establish
the first two assertions.  Pairwise nonparallelism within each family holds
almost surely.

It remains to establish \eqref{dir:good-charts}.  For a fixed unit vector
$v$, the one-dimensional Gaussian tail bound gives
\[
 \Pr\left\{\max_i|\langle v,g_{s,i}\rangle|
                      >10\sqrt{\log m}\right\}
 \le 2N e^{-50\log m}=2m^{-44}.
\]
Call such an index $s$ bad for $v$.  Before any conditioning, these bad-index
events are independent across $s$.  Therefore the probability that at least
$S/2$ indices are bad for this fixed $v$ is at most
\[
 2^{S}(2m^{-44})^{S/2}.
\]
Take a sphere net of radius $1/m$, with cardinality at most $(1+2m)^m$,
and union bound this last estimate over its points.  The resulting failure
probability tends to zero.  Finally intersect with the already established
norm event.  If $v_0$ is a net point with $|v-v_0|\le1/m$, then
\[
 |\langle v-v_0,g_{s,i}\rangle|\le\frac2{\sqrt m}
\]
for every $s,i$.  Every chart good for $v_0$ consequently satisfies
\eqref{dir:good-charts} for $v$, with $C=11$ for all sufficiently large
$m$.  This argument uses independence only for the unconditioned Gaussian
samples.  A final union bound intersects this event with the preceding
high-probability events and proves the lemma.
\end{proof}
